
de @@M0@@~» ;
pour définir la @@M1@@-ième classe de Chern, nous procèderons de façon analogue, en ramenant @@M2@@ au rang @@M3@@ en considérant @@M4@@, et nous prendrons la somme alternée des puissances extérieures de l'élément ainsi obtenu ; on «~interprétera~» la proposition qui suit en y faisant @@M5@@ :

\textbf{Proposition 3.2.} \uline{Soit} @@M6@@ \uline{un} @@M7@@-\uline{anneau}.
\begin{enumerate}[label=\roman*)]
\item @@M8@@
\item \uline{Si on pose} @@M9@@, \uline{et} @@M10@@, \uline{alors} :
@@M11@@
@@M12@@
\end{enumerate}

On voit donc que la connaissance des opérations @@M13@@ équivaut à celle des opérations @@M14@@.

\uline{Démonstration}.

i) On a
@@M15@@
De @@M16@@, on tire : @@M17@@, d'où : @@M18@@, ce qui donne la relation cherchée.

ii) @@M19@@
On développe @@M20@@. @@M21@@~; le coefficient de @@M22@@ est donc :
@@M23@@

